% !TEX root =  main.tex
\chapter{Summary and Conclusions}
\label{chap:5}
This thesis set out to answer a specific question left open by first-passage search theory: in a crowded, two-dimensional Brownian mixture, at what probe density does adding more searchers stop helping and start hindering detection, and can this transition be flagged earlier than by watching the mean detection time alone? Using Langevin dynamics simulations in LAMMPS, across six system sizes and a probe-density sweep from a single searcher to near-total population, this thesis confirms that a caldera-shaped trade-off exists, is not an artifact of any one system scale, and is quantifiable by more than one independent diagnostic.

The central results is that the optimal probe density, $\phi^*$, is consistent at roughly 29--31\% of the two-dimensional domain area across system sizes spanning a full order of magnitude ($N_\text{total} =30$ to 200), with a modestly elevated value at the smallest system studied ($N_\text{total}=10$). This optimum sits well below the theoretical hexagonal close-packing ceiling, which indicates that the breakdown of ideal, non-interacting search behavior is driven by dynamical caging effects during active search rather than by the searchers simply running out of physical room. The ergodic product, $E = N_p \times \langle T \rangle$, was confirmed as a more sensitive early indicator of this transition than the mean detection time alone, particularly for the larger, more heavily sampled system sizes, where it flagged the onset of crowding-influenced search tens of probes before the corresponding change became visually apparent in $\langle T \rangle$ itself. A threefold qualitative typology --- Incipient, Shield, and Steep-Walled caldera profiles --- was introduced to describe how the sharpness of this crowding response scales with the system size, and was shown to be corroborated independently by the ergodic product and the detection-time distribution shape, not solely by the parametric fit from which it was first drawn. In addition to the three primary caldera types, we observe a qualitatively distinct double-floor (bifurcated) shape in the $N=100$ system, with a weaker hint of the same pattern in $N=70$. We introduce this as a fourth caldera type, Type IV, and recommend it as a direction for future investigation with denser probe-number sampling.

% Beyond the three caldera types already identified, visual inspection of the $\langle T \rangle$ versus $N_p$ curves reveals a fourth, qualitatively distinct floor shape in the $N=100$ system where rather than a single continuous minimum, the curve traces two adjacent low points separated by a modest intervening rise. We classify this as a Bifurcated (Type IV) caldera. A similar, though considerably shallower, double dip shape is also visible in $N=70$, subtle enough that we regard it as a weaker instance of the same pattern rather than a clearly separate floor. The remaining systems each display conventional single-floor shape already captured by Types I-III.

These findings carry practical implications for point-of-care microfluidic and colloidal impurity-detection systems, where probe density is a design choice and not merely a fixed property of the sample. Rather than maximizing the number of functionalized probe particles introduced into a search volume, this work suggests that detection efficiency is optimized at a comparatively moderate, system-size-independent area fraction. Reducing not only detection time but also reagent cost and the risk of fluid clogging associated with overloading a sample with searchers. More broadly, by demonstrating a controlled, quantifiable breakdown of the ergodic $1/N_p$ scaling law well short of geometric packing limits, this work contributes a concrete data point to the wider study of non-equilibrium statistical mechanics in crowded, finite-sized environments which is a regime relevant to macromolecular crowding in biological cytoplasm as much as to synthetic colloidal suspensions.

This study has several limitations that bound the scope of these conclusions. The simulations are restricted to two dimensions, to purely repulsive Weeks-Chandler-Andersen interaction with no long-range electrostatics, hydrodynamics coupling, or gravitational effects, and to a fixed 2:1 probe-to-foreign diameter ratio. Each condition was averaged over ten independent seeds, and the production window was fixed 2.0\,s of simulated time, which may understate the true crowding penalty at the most heavily caged conditions if detection fractions there fall meaningfully below 1.0, as discussed in Section~\ref{sec:4.11}. The non-monotonic dependence of the absolute floor detection time on total system size (Section~\ref{sec:4.10}) also remains unresolved and is flagged as a direction for future work, alongside extending this framework to three dimensions, to a wider range of diameter ratios, and to explicit hydrodynamic interactions between probes.

Taken together, these results answer the question posed at the outset of this thesis: more searcher do not help without limit, and the point at which they stop helping is neither arbitrary  nor system-size-dependent in the way naive packing arguments would suggest. It is a stable, dynamically determined optimum, detectable earlier through the ergodic product than through raw detection time, and consistent enough across scales to serve as concrete, practical guidance for the design of crowded, particle-based detection systems.
